\chapter{DF-031 Status and feedback standard Blocks}\label{record-092}

\noindent{\footnotesize\color{BookMuted}\path{docs/design/DF-031-status-feedback-standard-blocks.md}\par}\smallskip

\begin{RecordMeta}
title: DF-031 Status and feedback standard Blocks
legend: DF
lane: design
priority: medium
keywords:
  - blocks
  - component-audit
  - dogfood
  - lowerings
  - status-feedback

\end{RecordMeta}

\section{Framing}\label{record-092-framing}

The v6 block floor already proved \texttt{A\allowbreak{}p\allowbreak{}p\allowbreak{}S\allowbreak{}h\allowbreak{}e\allowbreak{}l\allowbreak{}l}, \texttt{R\allowbreak{}e\allowbreak{}a\allowbreak{}d\allowbreak{}e\allowbreak{}r\allowbreak{}S\allowbreak{}u\allowbreak{}r\allowbreak{}f\allowbreak{}a\allowbreak{}c\allowbreak{}e}, and \texttt{I\allowbreak{}n\allowbreak{}s\allowbreak{}p\allowbreak{}e\allowbreak{}c\allowbreak{}t\allowbreak{}o\allowbreak{}r\allowbreak{}P\allowbreak{}a\allowbreak{}n\allowbreak{}e\allowbreak{}l}. The next useful slice is the product feedback layer: inline state, in-flow notices, transient overlays, activity history, shortcut cues, and progress.

This cycle turns issues \#220 through \#225 into first-party standard Blocks in \texttt{@\allowbreak{}f\allowbreak{}l\allowbreak{}y\allowbreak{}i\allowbreak{}n\allowbreak{}g\allowbreak{}r\allowbreak{}o\allowbreak{}b\allowbreak{}o\allowbreak{}t\allowbreak{}s\allowbreak{}/\allowbreak{}b\allowbreak{}i\allowbreak{}j\allowbreak{}o\allowbreak{}u}. The work uses the current typed Blocks API: \texttt{d\allowbreak{}e\allowbreak{}f\allowbreak{}i\allowbreak{}n\allowbreak{}e\allowbreak{}B\allowbreak{}l\allowbreak{}o\allowbreak{}c\allowbreak{}k\allowbreak{}(\allowbreak{})}, validated metadata, view data contracts, command intents, story ids, deterministic render output, and stable lowering facts. It does not introduce a new YAML or JSON runtime IR. Declarative YAML/JSON can become an adapter later, but it should compile into the typed contract this slice proves.

\section{Sponsored Users}\label{record-092-sponsored-users}

\begin{itemize}
\tightlist
\item
  TUI app authors who need reusable status and feedback semantics instead of bespoke strings inside pages.
\item
  DOGFOOD maintainers who need status surfaces to lower truthfully across interactive, static, pipe, and accessible modes.
\item
  Tooling authors who want block metadata, data contracts, command intents, and stories to describe feedback surfaces before rendering them.
\item
  Future IR authors who need a typed target for YAML/JSON block declarations.
\end{itemize}

\section{Hills}\label{record-092-hills}

\begin{enumerate}
\def\labelenumi{\arabic{enumi}.}
\tightlist
\item
  A builder can import one of the standard feedback Blocks and immediately inspect its metadata, data requirements, command intents, and ready story.
\item
  A reader can move between rich, static, pipe, and accessible profiles and still get the same feedback meaning, even when visual chrome disappears.
\item
  A maintainer can run one focused cycle test that proves the status and feedback slice is catalogued, rendered, and lowering-stable.
\end{enumerate}

\section{Playback Questions}\label{record-092-playback-questions}

\begin{itemize}
\tightlist
\item
  Which feedback surfaces are common enough to become first-party Blocks now?
\item
  Do these Blocks declare data contracts and command intents without taking on provider subscriptions or command dispatch?
\item
  Do visual modes render as surfaces while pipe and accessible modes preserve meaning as text?
\item
  Do lowerings carry stable facts so tests and tools do not parse terminal drawings?
\item
  Does this slice make room for a future YAML/JSON declaration format without inventing a second runtime contract?
\end{itemize}

\section{Block Set}\label{record-092-block-set}

{\def\LTcaptype{none} % do not increment counter
\begin{longtable}[]{@{}
  >{\raggedright\arraybackslash}p{(\linewidth - 4\tabcolsep) * \real{0.3333}}
  >{\raggedright\arraybackslash}p{(\linewidth - 4\tabcolsep) * \real{0.3333}}
  >{\raggedright\arraybackslash}p{(\linewidth - 4\tabcolsep) * \real{0.3333}}@{}}
\toprule\noalign{}
{I\allowbreak{}s\allowbreak{}s\allowbreak{}u\allowbreak{}e}
 & {B\allowbreak{}l\allowbreak{}o\allowbreak{}c\allowbreak{}k}
 & {P\allowbreak{}u\allowbreak{}r\allowbreak{}p\allowbreak{}o\allowbreak{}s\allowbreak{}e}
 \\
\midrule\noalign{}
\endhead
\bottomrule\noalign{}
\endlastfoot
\href{https://github.com/flyingrobots/bijou/issues/220}{{\#\allowbreak{}2\allowbreak{}2\allowbreak{}0}} & \texttt{I\allowbreak{}n\allowbreak{}l\allowbreak{}i\allowbreak{}n\allowbreak{}e\allowbreak{}S\allowbreak{}t\allowbreak{}a\allowbreak{}t\allowbreak{}u\allowbreak{}s\allowbreak{}B\allowbreak{}l\allowbreak{}o\allowbreak{}c\allowbreak{}k} & {S\allowbreak{}m\allowbreak{}a\allowbreak{}l\allowbreak{}l} {l\allowbreak{}o\allowbreak{}c\allowbreak{}a\allowbreak{}l} {s\allowbreak{}t\allowbreak{}a\allowbreak{}t\allowbreak{}u\allowbreak{}s} {b\allowbreak{}e\allowbreak{}s\allowbreak{}i\allowbreak{}d\allowbreak{}e} {a\allowbreak{}n} {i\allowbreak{}t\allowbreak{}e\allowbreak{}m\allowbreak{},} {f\allowbreak{}i\allowbreak{}e\allowbreak{}l\allowbreak{}d\allowbreak{},} {r\allowbreak{}o\allowbreak{}w\allowbreak{},} {o\allowbreak{}r} {c\allowbreak{}o\allowbreak{}m\allowbreak{}m\allowbreak{}a\allowbreak{}n\allowbreak{}d\allowbreak{}.} \\
\href{https://github.com/flyingrobots/bijou/issues/221}{{\#\allowbreak{}2\allowbreak{}2\allowbreak{}1}} & \texttt{I\allowbreak{}n\allowbreak{}F\allowbreak{}l\allowbreak{}o\allowbreak{}w\allowbreak{}S\allowbreak{}t\allowbreak{}a\allowbreak{}t\allowbreak{}u\allowbreak{}s\allowbreak{}B\allowbreak{}l\allowbreak{}o\allowbreak{}c\allowbreak{}k} & {N\allowbreak{}o\allowbreak{}t\allowbreak{}i\allowbreak{}c\allowbreak{}e} {o\allowbreak{}r} {s\allowbreak{}t\allowbreak{}a\allowbreak{}t\allowbreak{}e} {m\allowbreak{}e\allowbreak{}s\allowbreak{}s\allowbreak{}a\allowbreak{}g\allowbreak{}e} {t\allowbreak{}h\allowbreak{}a\allowbreak{}t} {l\allowbreak{}i\allowbreak{}v\allowbreak{}e\allowbreak{}s} {i\allowbreak{}n} {d\allowbreak{}o\allowbreak{}c\allowbreak{}u\allowbreak{}m\allowbreak{}e\allowbreak{}n\allowbreak{}t} {f\allowbreak{}l\allowbreak{}o\allowbreak{}w\allowbreak{}.} \\
\href{https://github.com/flyingrobots/bijou/issues/222}{{\#\allowbreak{}2\allowbreak{}2\allowbreak{}2}} & \texttt{T\allowbreak{}r\allowbreak{}a\allowbreak{}n\allowbreak{}s\allowbreak{}i\allowbreak{}e\allowbreak{}n\allowbreak{}t\allowbreak{}O\allowbreak{}v\allowbreak{}e\allowbreak{}r\allowbreak{}l\allowbreak{}a\allowbreak{}y\allowbreak{}B\allowbreak{}l\allowbreak{}o\allowbreak{}c\allowbreak{}k} & {T\allowbreak{}e\allowbreak{}m\allowbreak{}p\allowbreak{}o\allowbreak{}r\allowbreak{}a\allowbreak{}r\allowbreak{}y} {s\allowbreak{}h\allowbreak{}e\allowbreak{}l\allowbreak{}l\allowbreak{}-\allowbreak{}l\allowbreak{}e\allowbreak{}v\allowbreak{}e\allowbreak{}l} {f\allowbreak{}e\allowbreak{}e\allowbreak{}d\allowbreak{}b\allowbreak{}a\allowbreak{}c\allowbreak{}k} {t\allowbreak{}h\allowbreak{}a\allowbreak{}t} {s\allowbreak{}h\allowbreak{}o\allowbreak{}u\allowbreak{}l\allowbreak{}d} {s\allowbreak{}t\allowbreak{}i\allowbreak{}l\allowbreak{}l} {l\allowbreak{}o\allowbreak{}w\allowbreak{}e\allowbreak{}r} {t\allowbreak{}o} {t\allowbreak{}e\allowbreak{}x\allowbreak{}t\allowbreak{}.} \\
\href{https://github.com/flyingrobots/bijou/issues/223}{{\#\allowbreak{}2\allowbreak{}2\allowbreak{}3}} & \texttt{A\allowbreak{}c\allowbreak{}t\allowbreak{}i\allowbreak{}v\allowbreak{}i\allowbreak{}t\allowbreak{}y\allowbreak{}S\allowbreak{}t\allowbreak{}r\allowbreak{}e\allowbreak{}a\allowbreak{}m\allowbreak{}B\allowbreak{}l\allowbreak{}o\allowbreak{}c\allowbreak{}k} & {O\allowbreak{}r\allowbreak{}d\allowbreak{}e\allowbreak{}r\allowbreak{}e\allowbreak{}d} {r\allowbreak{}e\allowbreak{}c\allowbreak{}e\allowbreak{}n\allowbreak{}t} {e\allowbreak{}v\allowbreak{}e\allowbreak{}n\allowbreak{}t\allowbreak{}s} {w\allowbreak{}i\allowbreak{}t\allowbreak{}h} {o\allowbreak{}p\allowbreak{}t\allowbreak{}i\allowbreak{}o\allowbreak{}n\allowbreak{}a\allowbreak{}l} {a\allowbreak{}c\allowbreak{}t\allowbreak{}i\allowbreak{}v\allowbreak{}e} {s\allowbreak{}e\allowbreak{}l\allowbreak{}e\allowbreak{}c\allowbreak{}t\allowbreak{}i\allowbreak{}o\allowbreak{}n\allowbreak{}.} \\
\href{https://github.com/flyingrobots/bijou/issues/224}{{\#\allowbreak{}2\allowbreak{}2\allowbreak{}4}} & \texttt{S\allowbreak{}h\allowbreak{}o\allowbreak{}r\allowbreak{}t\allowbreak{}c\allowbreak{}u\allowbreak{}t\allowbreak{}C\allowbreak{}u\allowbreak{}e\allowbreak{}B\allowbreak{}l\allowbreak{}o\allowbreak{}c\allowbreak{}k} & {I\allowbreak{}n\allowbreak{}l\allowbreak{}i\allowbreak{}n\allowbreak{}e} {s\allowbreak{}h\allowbreak{}o\allowbreak{}r\allowbreak{}t\allowbreak{}c\allowbreak{}u\allowbreak{}t} {h\allowbreak{}i\allowbreak{}n\allowbreak{}t\allowbreak{}s} {s\allowbreak{}c\allowbreak{}o\allowbreak{}p\allowbreak{}e\allowbreak{}d} {t\allowbreak{}o} {a} {p\allowbreak{}a\allowbreak{}g\allowbreak{}e\allowbreak{},} {r\allowbreak{}e\allowbreak{}g\allowbreak{}i\allowbreak{}o\allowbreak{}n\allowbreak{},} {o\allowbreak{}r} {c\allowbreak{}o\allowbreak{}m\allowbreak{}m\allowbreak{}a\allowbreak{}n\allowbreak{}d\allowbreak{}.} \\
\href{https://github.com/flyingrobots/bijou/issues/225}{{\#\allowbreak{}2\allowbreak{}2\allowbreak{}5}} & \texttt{P\allowbreak{}r\allowbreak{}o\allowbreak{}g\allowbreak{}r\allowbreak{}e\allowbreak{}s\allowbreak{}s\allowbreak{}I\allowbreak{}n\allowbreak{}d\allowbreak{}i\allowbreak{}c\allowbreak{}a\allowbreak{}t\allowbreak{}o\allowbreak{}r\allowbreak{}B\allowbreak{}l\allowbreak{}o\allowbreak{}c\allowbreak{}k} & {D\allowbreak{}e\allowbreak{}t\allowbreak{}e\allowbreak{}r\allowbreak{}m\allowbreak{}i\allowbreak{}n\allowbreak{}a\allowbreak{}t\allowbreak{}e} {o\allowbreak{}r} {n\allowbreak{}a\allowbreak{}m\allowbreak{}e\allowbreak{}d} {p\allowbreak{}r\allowbreak{}o\allowbreak{}g\allowbreak{}r\allowbreak{}e\allowbreak{}s\allowbreak{}s} {s\allowbreak{}t\allowbreak{}a\allowbreak{}t\allowbreak{}e} {w\allowbreak{}i\allowbreak{}t\allowbreak{}h} {s\allowbreak{}t\allowbreak{}a\allowbreak{}b\allowbreak{}l\allowbreak{}e} {f\allowbreak{}a\allowbreak{}c\allowbreak{}t\allowbreak{}s\allowbreak{}.} \\
\end{longtable}
}

\section{TUI Mockups}\label{record-092-tui-mockups}

Inline status:

\begin{CodeBlock}
+-- InlineStatusBlock ----------------------+
| label: docs                               |
| status: ok                                |
| message: synced                           |
+-------------------------------------------+
\end{CodeBlock}

In-flow notice:

\begin{CodeBlock}
+-- InFlowStatusBlock ----------------------+
| severity: warning                         |
| source: docs                              |
| message: inventory stale                  |
| action: run docs:inventory                |
+-------------------------------------------+
\end{CodeBlock}

Transient overlay:

\begin{CodeBlock}
+-- TransientOverlayBlock ------------------+
| priority: normal                          |
| message: Saved DOGFOOD route              |
| dismiss: Esc dismisses                    |
+-------------------------------------------+
\end{CodeBlock}

Activity stream:

\begin{CodeBlock}
+-- ActivityStreamBlock --------------------+
| events: 10:41 tests passed                |
|         10:42 PR opened                   |
| selected: 10:41 tests passed              |
+-------------------------------------------+
\end{CodeBlock}

Shortcut cues:

\begin{CodeBlock}
+-- ShortcutCueBlock -----------------------+
| shortcuts: / Search                       |
|            ? Help                         |
|            Esc Close                      |
| scope: page                               |
+-------------------------------------------+
\end{CodeBlock}

Progress:

\begin{CodeBlock}
+-- ProgressIndicatorBlock -----------------+
| label: Install packages                   |
| value: 3                                  |
| total: 5                                  |
| percent: 60%                              |
+-------------------------------------------+
\end{CodeBlock}

\section{Lower Modes}\label{record-092-lower-modes}

Static mode preserves the same surface structure as interactive mode while remaining deterministic and render-only.

Pipe mode lowers every Block to a compact record:

\begin{CodeBlock}
InlineStatusBlock | label=docs | status=ok | message=synced
\end{CodeBlock}

Accessible mode lowers every Block to a readable sentence list:

\begin{CodeBlock}
InlineStatusBlock
- label: docs
- status: ok
- message: synced
\end{CodeBlock}

The lower-mode invariant is semantic parity, not visual parity. Pipe and accessible consumers should not need box drawing, row positions, or color to understand status.

\section{Requirements}\label{record-092-requirements}

\begin{itemize}
\tightlist
\item
  Export all six Blocks from the package barrel and standard catalog.
\item
  Include all six Blocks in \texttt{s\allowbreak{}t\allowbreak{}a\allowbreak{}n\allowbreak{}d\allowbreak{}a\allowbreak{}r\allowbreak{}d\allowbreak{}B\allowbreak{}l\allowbreak{}o\allowbreak{}c\allowbreak{}k\allowbreak{}P\allowbreak{}a\allowbreak{}c\allowbreak{}k\allowbreak{}a\allowbreak{}g\allowbreak{}e\allowbreak{}M\allowbreak{}a\allowbreak{}n\allowbreak{}i\allowbreak{}f\allowbreak{}e\allowbreak{}s\allowbreak{}t}.
\item
  Give each Block valid metadata across interactive, static, pipe, and accessible modes.
\item
  Give each Block at least one ready story id.
\item
  Attach view data contracts so tooling can inspect required slot data without provider handles.
\item
  Attach schema adapters and schema-bound block exports so callers can validate plain input before binding render slots.
\item
  Attach command intents for selecting the block, copying facts, and opening a story, without command callbacks.
\item
  Render interactive/static output as surfaces and pipe/accessible output as text.
\item
  Preserve \texttt{b\allowbreak{}l\allowbreak{}o\allowbreak{}c\allowbreak{}k}, \texttt{b\allowbreak{}l\allowbreak{}o\allowbreak{}c\allowbreak{}k\allowbreak{}.\allowbreak{}r\allowbreak{}e\allowbreak{}n\allowbreak{}d\allowbreak{}e\allowbreak{}r\allowbreak{}e\allowbreak{}d}, \texttt{b\allowbreak{}l\allowbreak{}o\allowbreak{}c\allowbreak{}k\allowbreak{}.\allowbreak{}f\allowbreak{}a\allowbreak{}m\allowbreak{}i\allowbreak{}l\allowbreak{}y}, \texttt{b\allowbreak{}l\allowbreak{}o\allowbreak{}c\allowbreak{}k\allowbreak{}.\allowbreak{}v\allowbreak{}a\allowbreak{}r\allowbreak{}i\allowbreak{}a\allowbreak{}n\allowbreak{}t}, \texttt{b\allowbreak{}l\allowbreak{}o\allowbreak{}c\allowbreak{}k\allowbreak{}.\allowbreak{}m\allowbreak{}o\allowbreak{}d\allowbreak{}e}, selected identity, and \texttt{s\allowbreak{}e\allowbreak{}m\allowbreak{}a\allowbreak{}n\allowbreak{}t\allowbreak{}i\allowbreak{}c\allowbreak{}V\allowbreak{}a\allowbreak{}l\allowbreak{}u\allowbreak{}e\allowbreak{}.\allowbreak{}*} facts across modes.
\end{itemize}

\section{Acceptance Criteria}\label{record-092-acceptance-criteria}

\begin{itemize}
\tightlist
\item
  \texttt{t\allowbreak{}e\allowbreak{}s\allowbreak{}t\allowbreak{}s\allowbreak{}/\allowbreak{}c\allowbreak{}y\allowbreak{}c\allowbreak{}l\allowbreak{}e\allowbreak{}s\allowbreak{}/\allowbreak{}D\allowbreak{}F\allowbreak{}-\allowbreak{}0\allowbreak{}3\allowbreak{}1\allowbreak{}/\allowbreak{}s\allowbreak{}t\allowbreak{}a\allowbreak{}t\allowbreak{}u\allowbreak{}s\allowbreak{}-\allowbreak{}f\allowbreak{}e\allowbreak{}e\allowbreak{}d\allowbreak{}b\allowbreak{}a\allowbreak{}c\allowbreak{}k\allowbreak{}-\allowbreak{}b\allowbreak{}l\allowbreak{}o\allowbreak{}c\allowbreak{}k\allowbreak{}s\allowbreak{}.\allowbreak{}t\allowbreak{}e\allowbreak{}s\allowbreak{}t\allowbreak{}.\allowbreak{}t\allowbreak{}s} proves the six Blocks publish through the standard catalog and package manifest.
\item
  The cycle test validates metadata and command intent shape for every Block.
\item
  The cycle test validates schema-bound input for every Block and rejects accessor-backed boundary objects before invoking getters.
\item
  The cycle test renders every Block in every canonical mode.
\item
  The cycle test verifies lowerings remain fact-stable and expose family, variant, selected, mode, and semantic-value facts.
\item
  Existing DX-031 standard block tests expand their catalog expectations from three Blocks to nine Blocks.
\item
  Repo docs stop describing this catalog expansion as future v7 work.
\end{itemize}

\section{Implementation Outline}\label{record-092-implementation-outline}

\begin{enumerate}
\def\labelenumi{\arabic{enumi}.}
\tightlist
\item
  Add a RED cycle test for issues \#220 through \#225.
\item
  Extend \texttt{B\allowbreak{}l\allowbreak{}o\allowbreak{}c\allowbreak{}k\allowbreak{}S\allowbreak{}c\allowbreak{}a\allowbreak{}l\allowbreak{}e} to support the product scales needed by this slice: \texttt{i\allowbreak{}n\allowbreak{}l\allowbreak{}i\allowbreak{}n\allowbreak{}e} and \texttt{o\allowbreak{}v\allowbreak{}e\allowbreak{}r\allowbreak{}l\allowbreak{}a\allowbreak{}y}.
\item
  Add six standard Block definitions in \texttt{s\allowbreak{}t\allowbreak{}a\allowbreak{}n\allowbreak{}d\allowbreak{}a\allowbreak{}r\allowbreak{}d\allowbreak{}-\allowbreak{}b\allowbreak{}l\allowbreak{}o\allowbreak{}c\allowbreak{}k\allowbreak{}s\allowbreak{}.\allowbreak{}t\allowbreak{}s}.
\item
  Register their stories, metadata, data contracts, schema adapters, command intents, and package manifest entries.
\item
  Reuse the existing section renderer pattern so visual and lower modes stay deterministic.
\item
  Update BEARING, ROADMAP, CHANGELOG, v6 backlog evidence, design-system docs, and package docs.
\end{enumerate}

\section{Future Declarative Shape}\label{record-092-future-declarative-shape}

The old Bijou IR idea still fits, but this slice keeps it as a declaration format over the typed API:

\begin{CodeBlock}
block: InlineStatusBlock
package: "@flyingrobots/bijou"
story: inline-status.ready
modes:
  - interactive
  - static
  - pipe
  - accessible
data:
  inlineStatus.label: docs
  inlineStatus.status: ok
  inlineStatus.message: synced
commands:
  - inlineStatus.select
  - inlineStatus.copyFacts
  - inlineStatus.openStory
facts:
  block: InlineStatusBlock
  block.family: status
  block.variant: ready
  semanticValue.label: docs
\end{CodeBlock}

A future declaration adapter can parse this shape, route it through \texttt{i\allowbreak{}n\allowbreak{}l\allowbreak{}i\allowbreak{}n\allowbreak{}e\allowbreak{}S\allowbreak{}t\allowbreak{}a\allowbreak{}t\allowbreak{}u\allowbreak{}s\allowbreak{}S\allowbreak{}c\allowbreak{}h\allowbreak{}e\allowbreak{}m\allowbreak{}a\allowbreak{}B\allowbreak{}l\allowbreak{}o\allowbreak{}c\allowbreak{}k}, and build typed \texttt{B\allowbreak{}l\allowbreak{}o\allowbreak{}c\allowbreak{}k\allowbreak{}D\allowbreak{}e\allowbreak{}f\allowbreak{}i\allowbreak{}n\allowbreak{}i\allowbreak{}t\allowbreak{}i\allowbreak{}o\allowbreak{}n} input or story fixtures. The runtime source of truth remains the typed Block contract.

\section{Drift Check}\label{record-092-drift-check}

The original DX-031 release floor was correct to limit itself to the first three Blocks. Drift appeared later when issues \#220 through \#225 became clear v6 candidates but ROADMAP and BEARING still left them in the v7 queue. This cycle corrects the tracker mirror and makes the expansion explicit, issue-backed v6 work.

No new follow-on debt is deferred by this slice. Broader component-family audits remain separate issues after \#225.

\section{Playback}\label{record-092-playback}

\begin{itemize}
\tightlist
\item
  RED: the status/feedback cycle test imported six standard Blocks that did not exist yet and failed at module load.
\item
  RED: the review follow-up required schema-bound feedback exports and stronger lowering facts; the cycle test failed before those symbols and facts existed.
\item
  GREEN: \texttt{@\allowbreak{}f\allowbreak{}l\allowbreak{}y\allowbreak{}i\allowbreak{}n\allowbreak{}g\allowbreak{}r\allowbreak{}o\allowbreak{}b\allowbreak{}o\allowbreak{}t\allowbreak{}s\allowbreak{}/\allowbreak{}b\allowbreak{}i\allowbreak{}j\allowbreak{}o\allowbreak{}u} now exports lower-camel block identifiers \texttt{i\allowbreak{}n\allowbreak{}l\allowbreak{}i\allowbreak{}n\allowbreak{}e\allowbreak{}S\allowbreak{}t\allowbreak{}a\allowbreak{}t\allowbreak{}u\allowbreak{}s\allowbreak{}B\allowbreak{}l\allowbreak{}o\allowbreak{}c\allowbreak{}k}, \texttt{i\allowbreak{}n\allowbreak{}F\allowbreak{}l\allowbreak{}o\allowbreak{}w\allowbreak{}S\allowbreak{}t\allowbreak{}a\allowbreak{}t\allowbreak{}u\allowbreak{}s\allowbreak{}B\allowbreak{}l\allowbreak{}o\allowbreak{}c\allowbreak{}k}, \texttt{t\allowbreak{}r\allowbreak{}a\allowbreak{}n\allowbreak{}s\allowbreak{}i\allowbreak{}e\allowbreak{}n\allowbreak{}t\allowbreak{}O\allowbreak{}v\allowbreak{}e\allowbreak{}r\allowbreak{}l\allowbreak{}a\allowbreak{}y\allowbreak{}B\allowbreak{}l\allowbreak{}o\allowbreak{}c\allowbreak{}k}, \texttt{a\allowbreak{}c\allowbreak{}t\allowbreak{}i\allowbreak{}v\allowbreak{}i\allowbreak{}t\allowbreak{}y\allowbreak{}S\allowbreak{}t\allowbreak{}r\allowbreak{}e\allowbreak{}a\allowbreak{}m\allowbreak{}B\allowbreak{}l\allowbreak{}o\allowbreak{}c\allowbreak{}k}, \texttt{s\allowbreak{}h\allowbreak{}o\allowbreak{}r\allowbreak{}t\allowbreak{}c\allowbreak{}u\allowbreak{}t\allowbreak{}C\allowbreak{}u\allowbreak{}e\allowbreak{}B\allowbreak{}l\allowbreak{}o\allowbreak{}c\allowbreak{}k}, and \texttt{p\allowbreak{}r\allowbreak{}o\allowbreak{}g\allowbreak{}r\allowbreak{}e\allowbreak{}s\allowbreak{}s\allowbreak{}I\allowbreak{}n\allowbreak{}d\allowbreak{}i\allowbreak{}c\allowbreak{}a\allowbreak{}t\allowbreak{}o\allowbreak{}r\allowbreak{}B\allowbreak{}l\allowbreak{}o\allowbreak{}c\allowbreak{}k}, plus matching \texttt{*\allowbreak{}S\allowbreak{}c\allowbreak{}h\allowbreak{}e\allowbreak{}m\allowbreak{}a\allowbreak{}A\allowbreak{}d\allowbreak{}a\allowbreak{}p\allowbreak{}t\allowbreak{}e\allowbreak{}r} and \texttt{*\allowbreak{}S\allowbreak{}c\allowbreak{}h\allowbreak{}e\allowbreak{}m\allowbreak{}a\allowbreak{}B\allowbreak{}l\allowbreak{}o\allowbreak{}c\allowbreak{}k} exports.
\item
  The six Blocks are in \texttt{s\allowbreak{}t\allowbreak{}a\allowbreak{}n\allowbreak{}d\allowbreak{}a\allowbreak{}r\allowbreak{}d\allowbreak{}B\allowbreak{}l\allowbreak{}o\allowbreak{}c\allowbreak{}k\allowbreak{}s}, \texttt{s\allowbreak{}t\allowbreak{}a\allowbreak{}n\allowbreak{}d\allowbreak{}a\allowbreak{}r\allowbreak{}d\allowbreak{}B\allowbreak{}l\allowbreak{}o\allowbreak{}c\allowbreak{}k\allowbreak{}S\allowbreak{}t\allowbreak{}o\allowbreak{}r\allowbreak{}i\allowbreak{}e\allowbreak{}s}, and \texttt{s\allowbreak{}t\allowbreak{}a\allowbreak{}n\allowbreak{}d\allowbreak{}a\allowbreak{}r\allowbreak{}d\allowbreak{}B\allowbreak{}l\allowbreak{}o\allowbreak{}c\allowbreak{}k\allowbreak{}P\allowbreak{}a\allowbreak{}c\allowbreak{}k\allowbreak{}a\allowbreak{}g\allowbreak{}e\allowbreak{}M\allowbreak{}a\allowbreak{}n\allowbreak{}i\allowbreak{}f\allowbreak{}e\allowbreak{}s\allowbreak{}t}.
\item
  Interactive/static modes produce surfaces; pipe/accessibility modes produce semantic text.
\item
  Stable lowering facts are present across all modes, including family, variant, selected identity where applicable, concrete mode, and semantic-value facts.
\end{itemize}

\section{Retrospective}\label{record-092-retrospective}

The useful move was not to invent a new IR under pressure. The useful move was to land the feedback vocabulary as typed, tested Blocks, then leave YAML/JSON as a future authoring adapter. That keeps the v6 release boundary concrete: geometry and Blocks continue to be inspectable package truth, not untyped story prose.

